
On
suppose toujours que @@M0@@. La multiplication pour la loi @@M1@@ par @@M2@@ définit un morphisme de foncteurs : @@M3@@, caractérisé par sa valeur pour l'élément @@M4@@ de @@M5@@. On a donc :
@@M6@@
@@M7@@
On en déduit :
@@M8@@
En effet, les deux membres sont des foncteurs additifs égaux pour @@M9@@. De (2.9.3) et 1.15, il résulte que le @@M10@@-ième polynôme universel définissant la multiplication par @@M11@@ est isobare de poids @@M12@@~; il sera noté @@M13@@.

Notons enfin que si @@M14@@ est une @@M15@@-@@M16@@-algèbre, i.e. un @@M17@@-anneau muni d'un @@M18@@-homomorphisme : @@M19@@, et si @@M20@@, on a :
@@M21@@
Ceci résulte en effet de @@M22@@, et de (2.9.4).

\section*{3. Les opérations @@M23@@}

On se place encore sur la catégorie @@M24@@.

3.1. Soit @@M25@@ un @@M26@@-anneau. Nous allons définir à partir des opérations @@M27@@ de @@M28@@ de nouvelles opérations, notées @@M29@@. Elles sont destinées à introduire une notion de classes de Chern sous des conditions qui seront indiquées plus tard. Pour le moment, donnons une «~justification~» des définitions qui suivent : si @@M30@@ est un fibré vectoriel de rang @@M31@@, la @@M32@@-ième classe de Chern de @@M33@@ devra être (quelle que soit la théorie des classes de Chern envisagée) la «~somme alternée des puissances extérieures 
